Modern industrial recommender systems operate over multi-million item catalogs under stringent latency Service Level Agreements (SLAs), typically demanding end-to-end response times below 40~ms. This paper presents the mathematical formulation, architectural design, and empirical validation of an enterprise-grade two-stage recommendation engine. The system integrates a Two-Tower Neural Retrieval Network (Stage~1) with a Deep \& Cross Network Version 2 (DCNv2) Multi-Task Ranking Engine (Stage~2), supplemented by a Disjoint Linear Upper Confidence Bound (LinUCB) Contextual Bandit for cold-start item exploration. In Stage~1, independent user and item towers project high-dimensional sparse identities and dense behavioral signals into a shared $\mathbb{R}^{64}$ embedding space, optimized via in-batch negative sampling and sampled softmax cross-entropy with $\log Q$ frequency correction. Precomputed item vectors are indexed into an Approximate Nearest Neighbor (ANN) index using Hierarchical Navigable Small World (HNSW) graphs, yielding sub-0.7~ms top-100 candidate retrieval. In Stage~2, candidates are scored using DCNv2, which leverages factorized low-rank explicit feature crossing in parallel with a deep feed-forward backbone to predict Click-Through Rate ($\text{pCTR}$) and Conversion Rate ($\text{pCVR}$). To mitigate recommendation starvation loops, a LinUCB bandit dynamically allocates a calibrated 5\% exploration traffic slice to zero-history items via Sherman-Morrison rank-1 recursive covariance updates. Empirical evaluations across 50,000 multi-modal interactions demonstrate high discriminative accuracy ($\text{Recall@50}=0.6315$, $\text{Recall@100}=0.7532$, $\text{NDCG@10}=0.9542$, $\text{CTR AUC}=0.9682$, and $\text{CVR AUC}=0.9563$). Latency profiling reveals an end-to-end P99 response time of 9.060~ms on commodity CPUs, comfortably meeting production SLAs.
